\section{Introduction}

Novel view synthesis aims to render unseen viewpoints of a scene from a limited number of images. Neural Radiance Fields (NeRF)~\cite{mildenhall2021nerf} achieve photorealistic quality but rely on dense multi-view supervision and long training times, which limits their use in practical settings such as robotics~\cite{irshad2022neuralfields} and agriculture~\cite{chen2022plantnerf} domains where capturing many views is often infeasible. 3DGS~\cite{kerbl20233d} addresses these computational bottlenecks, enabling real-time rendering and efficient optimization. However, in few-view scenarios, 3DGS tends to over-reconstruct available high-frequency (HF) regions, sharply reproducing textures and edges in training views, while losing smooth low-frequency (LF) structures and overall stability of the scene\cite{nguyen2025dwtnerf}. 

We address the few-view 3D reconstruction challenge with our proposed method, a frequency-aware extension of 3DGS that integrates the Discrete Wavelet Transform (DWT) to guide spatial-frequency learning. By decomposing rendered and ground-truth images into multi-scale sub-bands, the model explicitly balances both lowand high frequency information. Two complementary supervision strategies are introduced: a global DWT loss that preserves large-scale consistency and a patch-wise DWT loss that refines local details and fine edges. Together, these components balance between HF and LF regions, yielding sharper, more reliable reconstructions under sparse supervision. Building on this foundation, we extend the framework to handle multispectral data. The proposed MultiSpectral 3DGS jointly reconstructs RGB and Near-Infrared (NIR) views using shared geometry and modality-specific appearance parameters, ensuring spectral and spatial coherence across bands. To support research in this area, we introduce a new multispectral greenhouse dataset along with an open-source few-shot benchmarking package that standardizes sparse-view evaluation protocols.  We validate both the RGB and multispectral versions of LGDWT-GS on standard benchmarks (LLFF~\cite{mildenhall2019llff}, MipNeRF360~\cite{barron2022mipnerf360}) and on our new dataset. The main contributions of this work are summarized as follows:
\begin{itemize}
    \item Introduction of joint local and global frequency-domain supervision to improve 3D reconstruction.  
    \item Development of an open-source multispectral greenhouse dataset containing four spectral bands (Red, Green, Red-Edge, NIR).  
    \item Extension of 3DGS to the multispectral domain, enabling consistent cross-spectral reconstruction.  
    \item Establishment of a standardized few-shot 3DGS benchmark and evaluation protocol.  
\end{itemize}

